\documentclass[11pt]{article}
\usepackage{ap-psych}
\title{Unit 2: Cognition: Question Set}
\author{Paul Wang}
\date{November 2026}
\docpart{Unit 2 \textperiodcentered\ Question Set}

\begin{document}
\maketitle

\begin{abstract}
Exam-style multiple choice on all eight Unit 2 topics, with a shared-study set on memory and a full Article Analysis Question on eyewitness memory. Distractors are deliberately drawn from the look-alike terms in this unit (proactive vs.\ retroactive, availability vs.\ representativeness, reliability vs.\ validity). All items and the AAQ source are original and fictional, written for practice. Allow about 30 minutes for the multiple choice and 25 minutes for the AAQ.
\end{abstract}

\scoringline{MCQ: \blank\ / \qtotal \qquad AAQ: \blank\ / 7 \qquad Target: 21+ MCQ and 5+ AAQ \qquad Pace: 72 s per MCQ}

\section{Perception (2.1)}

\begin{q}{3}{2.1 Perceptual set}{Expectation from a label shapes what is perceived: perceptual set.}
Told that an ambiguous sketch shows ``a young woman,'' most participants see a young woman; told it shows ``an old woman,'' most see an old woman. This shows
\chh{bottom-up processing}{change blindness}{perceptual set}{retinal disparity}
\end{q}

\begin{q}{1}{2.1 Gestalt}{Grouping by nearness = proximity.}
Twelve dots arranged in three tight clusters of four are seen as ``three groups.'' Which Gestalt principle is at work?
\chh{Proximity}{Closure}{Similarity}{Figure--ground}
\end{q}

\begin{q}{4}{2.1 Depth cues}{Retinal disparity is binocular; it requires both eyes.}
A patient who has lost sight in one eye will have the most difficulty using which depth cue?
\chh{Interposition}{Texture gradient}{Linear perspective}{Retinal disparity}
\end{q}

\begin{q}{2}{2.1 Attention}{Not noticing a swapped person during a brief interruption = change blindness.}
A stranger asking for directions is swapped for a different person while a door passes between them; most people do not notice. This is
\ch{inattentional blindness}{change blindness}{the cocktail party effect}{perceptual constancy}
\end{q}

\section{Thinking and decisions (2.2)}

\begin{q}{1}{2.2 Representativeness}{Judging by resemblance to a stereotype while ignoring base rates.}
Told that a man is shy and loves poetry, people judge him more likely to be a classics professor than a truck driver, even though truck drivers vastly outnumber classics professors. This is the
\ch{representativeness heuristic}{availability heuristic}{sunk-cost fallacy}{framing effect}
\end{q}

\begin{q}{3}{2.2 Framing}{Same facts, different wording, different choice = framing.}
Shoppers prefer beef labeled ``80\% lean'' to beef labeled ``20\% fat.'' This demonstrates
\chh{priming}{mental set}{framing}{confirmation bias}
\end{q}

\begin{q}{2}{2.2 Gambler's fallacy}{Believing independent outcomes ``even out'' = gambler's fallacy.}
After a coin lands heads five times, Nia bets heavily on tails because ``it's due.'' Nia is showing
\chh{belief in an algorithm}{the gambler's fallacy}{functional fixedness}{divergent thinking}
\end{q}

\begin{q}{4}{2.2 Confirmation bias}{Seeking only supporting information = confirmation bias.}
A student who believes a supplement improves focus reads only reviews that praise it and skips critical studies. This is
\chh{priming}{the availability heuristic}{mental set}{confirmation bias}
\end{q}

\begin{q}{1}{2.2 Algorithms}{Trying every combination guarantees a solution: algorithm.}
Trying every possible combination on a three-digit lock until it opens is
\ch{an algorithm, which guarantees a solution but may be slow}{a heuristic, which is fast but error-prone}{an example of mental set}{an example of insight}
\end{q}

\section{Memory: structure and encoding (2.3--2.5)}

\begin{stimulus}
\textbf{Study C.} Researchers presented 40 nouns to 90 participants, randomly assigned to one of three conditions. Group 1 judged whether each word was in uppercase. Group 2 judged whether it rhymed with a target word. Group 3 judged whether it fit a sentence (``The \underline{\hspace{2em}} was on the table''). In a surprise recall test, mean words recalled were 6, 11, and 19.
\end{stimulus}

\begin{q}{2}{2.4 Levels of processing}{Recall rose with depth: structural < phonemic < semantic.}
Study C's results best support the
\ch{multi-store model's claim that sensory memory is brief}{levels-of-processing view that deeper, meaning-based encoding improves recall}{serial position effect}{spacing effect}
\end{q}

\begin{q}{4}{2.4 Encoding}{Rhyme judgments process sound: phonemic (shallow) encoding.}
In Study C, Group 2 used which type of encoding?
\chh{Semantic}{Structural}{Visual imagery}{Phonemic}
\end{q}

\begin{q}{3}{2.4 Methods}{Random assignment + manipulated encoding task = experiment.}
Study C can support a causal conclusion because
\ch{it used a surprise test}{the sample was large}{participants were randomly assigned to encoding conditions}{the words were all nouns}
\end{q}

\begin{q}{1}{2.3 Working memory}{Holding a visual route in mind = visuospatial sketchpad.}
Mentally rotating a map to plan a route mainly uses which component of working memory?
\chh{Visuospatial sketchpad}{Phonological loop}{Echoic memory}{Long-term potentiation}
\end{q}

\begin{q}{3}{2.4 Chunking}{Grouping digits into meaningful units = chunking, which expands STM capacity.}
Remembering 1-4-9-2-1-7-7-6 as ``1492'' and ``1776'' is an example of
\chh{the method of loci}{maintenance rehearsal}{chunking}{priming}
\end{q}

\begin{q}{2}{2.5 Sensory memory}{A brief visual trace lasting under a second = iconic memory.}
For a fraction of a second after a lightning flash, you can still ``see'' the scene. This is
\chh{echoic memory}{iconic memory}{working memory}{prospective memory}
\end{q}

\begin{q}{4}{2.3 LTP}{Strengthened synaptic efficiency with repeated stimulation = long-term potentiation.}
A neuroscientist finds that repeatedly stimulating a pathway in the hippocampus makes its synapses respond more strongly for weeks. This is
\chh{reuptake}{the refractory period}{consolidation failure}{long-term potentiation}
\end{q}

\section{Retrieval and forgetting (2.6--2.7)}

\begin{q}{1}{2.6 Recognition}{Recognition (multiple choice) is easier than recall (free response).}
Students usually score higher on multiple-choice than on fill-in-the-blank questions covering the same material because
\ch{recognition requires matching cues, which is easier than recall}{recall uses implicit memory}{multiple-choice items are encoded semantically}{recognition relies on sensory memory}
\end{q}

\begin{q}{3}{2.6 State-dependent}{Same internal (caffeinated) state at learning and test aids recall.}
Students who studied after drinking coffee recall more when tested after coffee than when tested without it. This is
\chh{context-dependent memory}{mood-congruent memory}{state-dependent memory}{the testing effect}
\end{q}

\begin{q}{2}{2.7 Interference}{New Spanish learning disrupts old French recall = retroactive.}
After a semester of Spanish, Kai struggles to recall the French vocabulary he learned the year before. This is
\chh{proactive interference}{retroactive interference}{source amnesia}{encoding failure}
\end{q}

\begin{q}{4}{2.7 Forgetting curve}{Ebbinghaus: steep early loss, then leveling off.}
According to Ebbinghaus's forgetting curve, most forgetting of new, meaningless material occurs
\ch{at a steady rate over weeks}{only after several months}{only if the material was encoded deeply}{soon after learning, then levels off}
\end{q}

\begin{q}{1}{2.7 Source amnesia}{Remembering content but misattributing its source = source amnesia.}
Rina confidently tells a friend a fact she ``read in a textbook'' that she actually heard in a movie. This is
\chh{source amnesia}{retrograde amnesia}{repression}{the tip-of-the-tongue effect}
\end{q}

\section{Intelligence (2.8)}

\begin{q}{3}{2.8 Validity}{A test that forecasts the later outcome it targets has predictive validity.}
A college admissions test is valued because its scores correlate with students' first-year GPA. This reflects the test's
\chh{reliability}{standardization}{predictive validity}{split-half consistency}
\end{q}

\begin{q}{2}{2.8 Achievement vs. aptitude}{Measuring what has been learned = achievement test.}
The AP Psychology exam is best described as
\ch{an aptitude test, because it predicts college success}{an achievement test, because it measures learned content}{a measure of fluid intelligence}{a projective test}
\end{q}

\begin{q}{4}{2.8 Normal curve}{Deviation IQ: mean 100, SD 15; 130 = +2 SD \(\approx\) top 2--3\%.}
On the Wechsler scales, an IQ of 130 places a person about
\ch{at the 50th percentile}{1 SD above the mean}{in the bottom 2\%}{2 SDs above the mean, higher than about 97--98\% of people}
\end{q}

\begin{q}{1}{2.8 Flynn effect}{Rapid generational IQ gains point to environmental causes.}
Average raw IQ-test performance rose about 3 points per decade during the 20th century. This Flynn effect most strongly suggests that
\ch{environmental factors such as nutrition and schooling influence test performance}{intelligence is entirely genetic}{IQ tests have no reliability}{fluid intelligence declines with age}
\end{q}

\section{Article Analysis Question}

\begin{frq}{Question 1 \textperiodcentered\ AAQ \textperiodcentered\ 7 points \textperiodcentered\ 25 minutes}
\small
\textbf{Source (summarized; fictional, for practice).} Haddad, R., \& Kim, J. (2023). Post-event questioning and eyewitness accuracy. \emph{Journal of Memory and Law}, 9, 44--61.

\textbf{Participants.} 150 adult volunteers (ages 18--64) recruited online in Canada. All gave informed consent; because the study involved mild deception about its purpose, participants were fully debriefed at the end.

\textbf{Method.} All participants watched the same 90-second video of a store theft. They were then randomly assigned to read one of two written summaries. The \emph{misleading} summary stated that the thief wore a red cap (in the video, the cap was blue). The \emph{neutral} summary mentioned ``a cap'' with no color. Thirty minutes later, all participants answered a 12-item memory test; accuracy was the number of the 12 details reported correctly. Confidence was rated on a 1--7 scale.

\textbf{Results.} 47\% of the misleading group reported a red cap, compared with 6\% of the neutral group. Mean accuracy was 7.8 (misleading) vs.\ 9.4 (neutral), \(p < .01\). Confidence did not differ between groups (\(M = 5.6\) vs.\ \(5.5\)).

\medskip
\begin{enumerate}[label=(\Alph*),leftmargin=2em,itemsep=3pt]
\item Identify the research method used in the study.
\item State the operational definition of eyewitness accuracy.
\item Describe what the confidence results indicate.
\item Identify at least one ethical guideline applied by the researchers.
\item Explain the extent to which the findings may or may not be generalizable, using specific information from the study.
\item Explain how at least one finding supports or refutes the claim that ``memories are reconstructed rather than replayed exactly as they were stored.''
\end{enumerate}
\end{frq}

\makekey

\subsection*{AAQ scoring guide (7 points)}
\small
\begin{compare}{Part}{Earns the point}{Does not earn the point}
A (1) & ``Experiment.'' & ``Case study''; ``survey.'' \\
B (1) & The number of the 12 details on the memory test reported correctly, 30 minutes after the video. & ``How well they remembered the video.'' \\
C (1) & Participants given misleading information were just as confident (about 5.6 of 7) as the neutral group, even though they were less accurate: confidence did not track accuracy. & ``The misleading group was more confident''; restating the numbers with no meaning. \\
D (1) & Informed consent \emph{or} debriefing after deception, linked to the study. & ``They were ethical''; naming confidentiality, which the source does not mention. \\
E (1) & Limited: online Canadian volunteers watching a 90-second video in a low-stakes setting may differ from real eyewitnesses under stress. \emph{Or} broader: a wide adult age range (18--64) and random assignment support generalizing to adults in similar settings. & ``Small sample'' with no specifics. \\
F (2) & 1 pt: supports the claim, citing a result (47\% vs.\ 6\% reported a red cap; accuracy 7.8 vs.\ 9.4). 1 pt: explains with a concept, e.g., the \emph{misinformation effect}/\emph{constructive memory}: information encountered after the event was blended into the memory at retrieval, so the memory was rebuilt rather than replayed. & Defining the misinformation effect without linking it to the cap result. \\
\end{compare}

\end{document}
